\documentclass[10pt,a4paper]{amsart}
\usepackage[margin=19mm]{geometry}
\usepackage{amsmath,amssymb,mathtools}
\usepackage[scr=rsfs,cal=euler]{mathalfa}
\usepackage{xfrac}
\usepackage[unicode,hidelinks,bookmarks=false]{hyperref}
\newcommand{\ie}{i.e.\ }
\newcommand{\cf}{cf.\ }
\newcommand{\resp}{respectively\ }
\newcommand{\Subsection}{\subsection}
\providecommand{\nobreakdash}{\nobreak\hbox{-}}
\setlength{\emergencystretch}{2em}
\multlinegap0pt
\allowdisplaybreaks
\usepackage{indentfirst}
\usepackage{fullpage}
\usepackage{cleveref}
\setlength{\footskip}{24pt}
\numberwithin{equation}{section}
\newcommand{\ddbar}{\sqrt{-1}\,\partial\overline\partial}
\newtheorem{theorem}{Theorem}
\newtheorem{propo}{Proposition}
\newtheorem{lema}{Lemma}
\newtheorem{remark}{Remark}
\newtheorem{example}{Example}
\begin{document}
\title[Chern-Ricci flow on Kato surfaces]{Chern-Ricci flow on Kato surfaces}
\author{Daniele Angella; Mauricio Corrêa}
\date{}
\maketitle
\noindent\emph{Source and licence.} Chern-Ricci flow on Kato surfaces. Version arXiv:2609.18579v1 (CC BY 4.0). This material is distributed under the Creative Commons Attribution 4.0 International licence, http://creativecommons.org/licenses/by/4.0/, as recorded in the arXiv OAI licence metadata for this exact version and shown on the abstract page https://arxiv.org/abs/2609.18579v1. Full rights notice: licenses/arxiv-2609.18579-CC-BY-4.0.txt.\par\smallskip
\noindent\textit{Editorial scope.} Counted unit: the original \texttt{theorem} environment \texttt{thm:main-open-flow} at source lines 726--744 together with its complete proof at lines 746--856; the delivered document keeps source lines 432--857 so that every referenced label and every citation resolves inside it. Native editable source is the paper's own concatenated TeX; the AMS wrapper is editorial formatting only. No publisher class, upstream script or unrelated artwork is redistributed.
\section{Kato surfaces of intermediate type}

Let $S$ be a Kato surface of intermediate type and let $D$ be its maximal reduced divisor of rational curves. For surfaces of index one, Dloussky--Oeljeklaus--Toma \cite{dloussky-oeljeklaus-toma} identify the universal cover of $S\setminus D$ with $X=\mathbb H_\ell\times\mathbb C$, where $\mathbb H_\ell=\{\Re w<0\}$, and put the deck action into an affine triangular normal form. The higher-index case is obtained from this model by a finite cyclic extension of the deck group. Set $r:=\operatorname{ind}(S)$. We recall the form needed below.

 

 \begin{propo}\label{prop:higher-index-green-model}
There exists an intermediate surface $S_1$ of index one, with maximal reduced divisor of rational curves $D_1$, and an unramified cyclic covering
$$
A_1:=S_1\setminus D_1\longrightarrow A:=S\setminus D
$$
of degree $r$. On their common universal cover $X=\mathbb H_\ell\times\mathbb C$ one may choose coordinates so that the deck group $\Gamma_1$ of $A_1$ is generated by
$$
T(w,z)=(w+2\pi\sqrt{-1},z),\qquad
G(w,z)=(kw,\lambda z+h(w)),
$$
where $k=k(S)\in\mathbb Z$, $k\geq2$, and the full deck group $\Gamma$ of $A$ is generated by $\Gamma_1$ and
$$
\varphi(w,z)=\left(w-\frac{2\pi\sqrt{-1}}{r},\xi z\right),
\qquad
\xi:=e^{2\pi\sqrt{-1}\sigma_0/r},
$$
where $\sigma_0$ is the largest exponent occurring in the polynomial $Q$ of the associated index-one normal form; see \cite[Sections~3.1--3.2]{oeljeklaus-toma-moduli}. In particular, $\xi^r=1$ and $|\xi|=1$.
Here $r$ divides $k-1$, and
$$
\varphi^r=T^{-1},\ 
\varphi T\varphi^{-1}=T,\ 
\varphi G\varphi^{-1}=T^{(k-1)/r}G.
$$
In particular $G\varphi G^{-1}=\varphi^k$, so that
$\Gamma\simeq\mathbb Z[1/k]\rtimes\mathbb Z$. Moreover, $\varphi$ preserves
$s:=-\Re w$.
\end{propo}

\begin{proof}
For $r=1$ this is the affine Green normal form of
\cite{dloussky-oeljeklaus-toma}. For $r>1$, Oeljeklaus--Toma
\cite[Section~3.2 and Remark~4.4]{oeljeklaus-toma-moduli} show that every
intermediate surface of index $r$ is obtained as a cyclic quotient of an
index-one model. In their coordinates $(z,w)\in\mathbb C\times\mathbb H_\ell$,
the additional normalising transformation has the form
$$
(z,w)\longmapsto
\left(e^{2\pi\sqrt{-1}\sigma_0/r}z,
w-\frac{2\pi\sqrt{-1}}{r}\right).
$$
Interchanging the two coordinates gives the form above. 
The normal form satisfies the property $(I_q)$ of
\cite[Section~3.2, p.~329]{oeljeklaus-toma-moduli}, with $q=r$ in our notation. Thus, writing
$$
h(w)=Q(e^{-w}),\qquad Q(\zeta)=\sum_{m=\ell}^{\sigma} b_m\zeta^m,
$$
one has $a_0=0$ and
$
r\mid(k-1),\ r\mid(m-m')
$
whenever $b_m b_{m'}\neq0$. Since $\sigma_0$ is an exponent with $b_{\sigma_0}\neq0$, every exponent $m$ with $b_m\neq0$ satisfies
$
m\equiv\sigma_0\pmod r.
$
Therefore, with $\xi=e^{2\pi\sqrt{-1}\sigma_0/r}$,
$$
h\left(w-\frac{2\pi\sqrt{-1}}{r}\right)
=
\sum_m b_m e^{-mw}e^{2\pi\sqrt{-1}m/r}
=
\xi\,h(w).
$$
It follows that
$
\varphi G\varphi^{-1}=T^{(k-1)/r}G.
$
Since $\xi^r=1$, one also has $\varphi^r=T^{-1}$. Therefore
$$
\varphi G\varphi^{-1}=\varphi^{-(k-1)}G,
$$
equivalently
$$
G\varphi G^{-1}=\varphi^k.
$$
The resulting presentation is the standard presentation of
$\mathbb Z[1/k]\rtimes\mathbb Z$. Finally, $\varphi^*s=s$ is immediate.
\end{proof}

We work henceforth on this common universal cover and with the full deck group $\Gamma$. The construction below produces a $\Gamma$-equivariant Hermitian metric and an explicit solution of the normalised Chern--Ricci flow on $S\setminus D$.

\subsection{The affine Green model}
Each element $\gamma\in\Gamma$ has the affine triangular form
\begin{equation}\label{eq:general-affine-element}
 \gamma(w,z)=\bigl(a_\gamma w+b_\gamma,\lambda_\gamma z+h_\gamma(w)\bigr),
\end{equation}
where $a_\gamma>0$, $\lambda_\gamma\in\mathbb{C}^*$, and $h_\gamma \colon \mathbb H_\ell \to \mathbb C$ is holomorphic.
For the distinguished elements of the index-one subgroup we write
$$ T(w,z)=(w+2\pi\sqrt{-1},z),
\quad \text{ and } \quad
G(w,z)=(kw,\lambda z+h(w)), $$
where $k=k(S)\in\mathbb Z$, $k\geq2$. If $r>1$, the full group also contains the transformation $\varphi$ of \Cref{prop:higher-index-green-model}.
Set $s:=-\Re w>0$ and $\theta:=\Im w$, so that $w=-s+\sqrt{-1}\theta$.
Then
\begin{equation}\label{eq:s-transforms}
 T^*s=s,
 \qquad
 G^*s=ks.
\end{equation}
Define the transverse Green form by
\begin{equation*}
 \omega_{\rm tr}:=\ddbar(-\log s).
\end{equation*}
Since $s=-(w+\bar w)/2$, we have $\partial s=-\tfrac12\,dw$ and $\overline\partial s=-\tfrac12\,d\bar w$. Hence
$$ \omega_{\rm tr} = \sqrt{-1} \partial\overline\partial (-\log s) =
 \frac{\sqrt{-1}}{4s^2}\,dw\wedge d\bar w. $$
The transformation law of $s$ gives
\begin{equation*}
 T^*\omega_{\rm tr}=\omega_{\rm tr},
 \quad
 G^*\omega_{\rm tr}
 =
 \ddbar(-\log(ks))
 =
 \ddbar(-\log s)
 =
 \omega_{\rm tr}.
\end{equation*}
For $r>1$, one also has $\varphi^*s=s$ and hence $\varphi^*\omega_{\rm tr}=\omega_{\rm tr}$. Thus $\omega_{\rm tr}$ is invariant under the full deck group and descends to $S\setminus D$.


\subsection{The cut-off and the equivariant vertical coform}
The vertical tangent distribution is $V:=\mathbb{C}\cdot\partial_z\subset T_X^{1,0}$.
It is preserved by $\Gamma$: for an element
$\gamma\in\Gamma$ as in \eqref{eq:general-affine-element}, one has
\begin{equation*}
 d\gamma(\partial_z)=\lambda_\gamma\partial_z.
\end{equation*}
In particular, $dT(\partial_z)=\partial_z$ and $dG(\partial_z)=\lambda\partial_z$; when $r>1$, $d\varphi(\partial_z)=\xi\partial_z$.

Since the action of $\Gamma$ on $X$ is free and properly discontinuous, the
quotient map $q:X\to X/\Gamma=A=S\setminus D$ is a smooth covering. We fix once and for all a standard cut-off function for this covering: a smooth function
$\chi\in\mathcal C^\infty(X)$, $\chi\geq0$, such that the family
$\{\chi\circ\gamma\}_{\gamma\in\Gamma}$ is locally finite and
\begin{equation}\label{eq:cutoff-sum}
 \sum_{\gamma\in\Gamma} \chi(\gamma p)=1
\end{equation}
for every $p\in X$. Such a function is obtained in the usual way by lifting a locally
finite partition of unity on the quotient to chosen sheets over evenly covered
coordinate neighbourhoods and extending by zero; proper discontinuity gives
the required local finiteness.

Starting with the Euclidean Hermitian metric $g_0=|dw|^2+|dz|^2$, we set
\begin{equation*}
 g:=\sum_{\gamma\in\Gamma}\chi(\gamma p)\,\gamma^*g_0.
\end{equation*}
The sum is locally finite, hence $g$ is a smooth Hermitian metric. It is positive because the coefficients $\chi(\gamma p)$ are non-negative, at least
one of them is positive at each point, and every $\gamma^*g_0$ is positive.
It is also $\Gamma$-invariant: for $\eta\in\Gamma$,
\begin{align*}
 (\eta^*g)_p
 &=
 \sum_{\gamma\in\Gamma}\chi(\gamma\eta p)(\gamma\eta)^*g_0|_p \\
 &=
 \sum_{\delta\in\Gamma}\chi(\delta p)\delta^*g_0|_p
 =
 g_p,
\end{align*}
where $\delta=\gamma\eta$.

Let $H\subset T_X^{1,0}$ be the $g$-orthogonal complement of $V$. Since
both $g$ and $V$ are $\Gamma$-invariant, so is $H$. There is a unique
smooth $(1,0)$-form $\Theta$ such that $\Theta(\partial_z)=1$ and $\Theta|_H=0$. Write
\begin{equation*}
 \Theta=dz+B\,dw.
\end{equation*}
In local coordinates, write
\begin{equation*}
 g
 =
 g_{w\bar w}\,dw\,d\bar w
 +
 g_{w\bar z}\,dw\,d\bar z
 +
 g_{z\bar w}\,dz\,d\bar w
 +
 g_{z\bar z}\,dz\,d\bar z.
\end{equation*}
A horizontal vector has the form $\partial_w+\alpha\partial_z$. Orthogonality
to $\partial_z$ gives $g_{w\bar z}+\alpha g_{z\bar z}=0$, hence
$$ \alpha=-\frac{g_{w\bar z}}{g_{z\bar z}}. $$
Since $\Theta(\partial_w+\alpha\partial_z)=B+\alpha$, the condition $\Theta|_H=0$ gives
\begin{equation*}
 B=\frac{g_{w\bar z}}{g_{z\bar z}}.
\end{equation*}

For an affine element $\gamma$ as in \eqref{eq:general-affine-element},
\begin{equation*}
 d(a_\gamma w+b_\gamma)=a_\gamma\,dw,
 \qquad
 d(\lambda_\gamma z+h_\gamma(w))
 =
 \lambda_\gamma\,dz+h_\gamma'(w)\,dw.
\end{equation*}
Thus
$$
\gamma^*g_0=|a_\gamma|^2|dw|^2+|\lambda_\gamma dz+h_\gamma'(w)dw|^2.
$$
Expanding the second term gives
\begin{align*}
\gamma^* g_0 &=
|a_\gamma|^2|dw|^2 + |\lambda_\gamma dz+h_\gamma'(w)dw|^2 \\
&=
|a_\gamma|^2|dw|^2 + |\lambda_\gamma|^2 dz\,d\bar z
+h_\gamma'(w)\overline{\lambda_\gamma}\,dw\,d\bar z
+\lambda_\gamma\overline{h_\gamma'(w)}\,dz\,d\bar w
+|h_\gamma'(w)|^2dw\,d\bar w.
\end{align*}
Hence
\begin{equation*}
 g_{z\bar z}
 =
 \sum_{\gamma\in\Gamma}\chi(\gamma p)|\lambda_\gamma|^2,
 \qquad
 g_{w\bar z}
 =
 \sum_{\gamma\in\Gamma}\chi(\gamma p)h_\gamma'(w)\overline{\lambda_\gamma}.
\end{equation*}
Consequently,
\begin{equation*}
 B(w,z,\bar w,\bar z)
 =
 \frac{
 \displaystyle\sum_{\gamma\in\Gamma}
 \chi(\gamma(w,z))h_\gamma'(w)\overline{\lambda_\gamma}
 }{
 \displaystyle\sum_{\gamma\in\Gamma}
 \chi(\gamma(w,z))|\lambda_\gamma|^2
 }.
\end{equation*}
The denominator is strictly positive by \eqref{eq:cutoff-sum}, since $\lambda_\gamma\ne0$. Hence $B$ is smooth; in general it is neither holomorphic nor a function of $w$ alone.

\begin{lema} \label{lem:Theta-equivariant}
For every $\gamma\in\Gamma$, one has $\gamma^*\Theta=\lambda_\gamma\Theta$. In particular, $T^*\Theta=\Theta$, $G^*\Theta=\lambda\Theta$, and $\varphi^*\Theta=\xi\Theta$ when $r>1$.
\end{lema}

\begin{proof}
The kernel of $\Theta$ is $H$. Since $H$ is $\Gamma$-invariant,
$\gamma^*\Theta$ has the same kernel as $\Theta$. Hence it is a scalar
multiple of $\Theta$. Evaluating on $\partial_z$, we find
\begin{equation*}
 (\gamma^*\Theta)(\partial_z)
 =
 \Theta(d\gamma(\partial_z))
 =
 \Theta(\lambda_\gamma\partial_z)
 =
 \lambda_\gamma,
\end{equation*}
proving the statement.
\end{proof}

\subsection{The explicit normalised flow on the open surface}

Set
\begin{equation}\label{eq:mu-def}
 \mu:=-\frac{2\log|\lambda|}{\log k}.
\end{equation}
For $r>1$, replacing $G$ by $\varphi^jG$ changes its vertical multiplier from $\lambda$ to $\xi^j\lambda$; hence $|\lambda|$, and therefore $\mu$, is independent of this choice. By definition, $k^\mu|\lambda|^2=1$.
By Lemma~\ref{lem:Theta-equivariant} and \eqref{eq:s-transforms},
\begin{align*}
 G^*\left(s^\mu\sqrt{-1}\,\Theta\wedge\overline\Theta\right)
 &=
 (ks)^\mu\sqrt{-1}\,\lambda\Theta\wedge
 \bar\lambda\,\overline\Theta \\
 &=
 k^\mu|\lambda|^2s^\mu
 \sqrt{-1}\,\Theta\wedge\overline\Theta \\
 &=
 s^\mu\sqrt{-1}\,\Theta\wedge\overline\Theta.
\end{align*}
The same form is $T$-invariant. If $r>1$, then $\varphi^*s=s$ and $\varphi^*\Theta=\xi\Theta$ with $|\xi|=1$, so it is also $\varphi$-invariant. Hence it is invariant under the full deck group $\Gamma$ and descends to $S\setminus D$.
Assume from now on that $0<\mu<2$. For $M>0$, define
$$
A(t):=2-\mu+e^{-t}\bigl(M-(2-\mu)\bigr)=e^{-t}M+(1-e^{-t})(2-\mu).
$$
Thus $A(t)>0$ for all $t\geq0$, because $M>0$ and $2-\mu>0$.
Define the time-dependent $(1,1)$-form
\begin{equation}\label{eq:Omega-flow-def}
 \Omega_t^M
 :=
 A(t)\omega_{\rm tr}
 +
 e^{-t}s^\mu\sqrt{-1}\,\Theta\wedge\overline\Theta.
\end{equation}
The preceding invariance calculation shows that $\Omega_t^M$ descends from
$X$ to $S\setminus D$.



\begin{theorem}\label{thm:main-open-flow}
Let $S$ be a Kato surface of intermediate type, let $D$ be its maximal reduced divisor of rational curves, and assume $0<\mu<2$,  namely $k^{-1}<|\lambda|<1$.
Let $\Theta=dz+B\,dw$ be any
smooth $\Gamma$-equivariant vertical $(1,0)$-form on $X$. 
For every $M>0$, the form $\Omega_t^M$ on $S\setminus D$ defined in \eqref{eq:Omega-flow-def} is a smooth Hermitian metric on $S\setminus D$. It solves the normalised Chern--Ricci flow
\begin{equation*}
 \partial_t\Omega_t^M=-\operatorname{Ric}(\Omega_t^M)-\Omega_t^M
\end{equation*}
with initial value $\Omega_0^M = M\omega_{\rm tr} + s^\mu\sqrt{-1}\,\Theta\wedge\overline\Theta$.
Furthermore,
\begin{equation*}
 \Omega_t^M
 \longrightarrow
 (2-\mu)\omega_{\rm tr}
 \quad
 \text{in } \mathcal C^\infty_{\rm loc}(S\setminus D).
\end{equation*}
The limiting form is semipositive of rank one, and its kernel is the canonical foliation $\mathcal F$ considered above.
\end{theorem}

\begin{proof}
By construction, $\Omega_t^M$ is smooth and $\Gamma$-invariant on $X$;
hence it descends to a smooth $(1,1)$-form on $S\setminus D$. To prove positivity, let $v\in T_X^{1,0}$. Then
\begin{equation*}
 -\sqrt{-1}\,\Omega_t^M(v,\bar v)
 =
 A(t)\frac{|dw(v)|^2}{4s^2}
 +
 e^{-t}s^\mu|\Theta(v)|^2.
\end{equation*}
Both coefficients on the right-hand side are positive. If the right-hand side vanishes, then
$dw(v)=0$ and $\Theta(v)=0$. The first equality says that $v$ is
vertical, so $v=\alpha\partial_z$; the second gives
$\alpha=\Theta(v)=0$. Thus $v=0$, and $\Omega_t^M$ is positive definite.

We compute the determinant. Writing $\Theta=dz+Bdw$, one has
$$
\sqrt{-1}\,\Theta\wedge\overline\Theta=\sqrt{-1}\,dz\wedge d\bar z+\sqrt{-1}\,\bar B\,dz\wedge d\bar w+\sqrt{-1}\,B\,dw\wedge d\bar z+\sqrt{-1}\,|B|^2dw\wedge d\bar w.
$$
Hence the Hermitian matrix of $\Omega_t^M$ in the coframe $(dw,dz)$ is
\begin{equation*}
 \begin{pmatrix}
 \displaystyle \frac{A(t)}{4s^2}+e^{-t}s^\mu |B|^2
 &
 \displaystyle e^{-t}s^\mu B
 \\
 \displaystyle e^{-t}s^\mu\bar B
 &
 \displaystyle e^{-t}s^\mu
 \end{pmatrix}.
\end{equation*}
Hence
\begin{align*}
 \det(\Omega_t^M)
 &=
 \left(
 \frac{A(t)}{4s^2}
 +
 e^{-t}s^\mu|B|^2
\right)e^{-t}s^\mu
 -
 e^{-2t}s^{2\mu}|B|^2 \\
 &=
 \frac{A(t)e^{-t}}4s^{\mu-2}.
\end{align*}
The $B$-terms cancel exactly.
The Chern--Ricci form is
\begin{align*}
\operatorname{Ric}(\Omega_t^M)
 &= -\ddbar\log\det(\Omega_t^M) \\
 &= -\sqrt{-1}\partial\overline\partial \left( \log\left(\frac{A(t)e^{-t}}4\right) + (\mu-2)\log s\right) \\
 &= -\ddbar\bigl((\mu-2)\log s\bigr) \\
 &= -(\mu-2)\ddbar\log s \\
 &= -(2-\mu)\omega_{\rm tr}.
\end{align*}

We verify the evolution equation. From the definition of $A(t)$,
\begin{equation*}
 A'(t)
 =
 -e^{-t}\bigl(M-(2-\mu)\bigr)
 =
 (2-\mu)-A(t).
\end{equation*}
Since $\omega_{\rm tr}$, $s$, and $\Theta$ are independent of $t$,
\begin{equation*}
 \partial_t\Omega_t^M
 =
 A'(t)\omega_{\rm tr}
 -
 e^{-t}s^\mu\sqrt{-1}\,\Theta\wedge\overline\Theta.
\end{equation*}
On the other hand,
\begin{align*}
 -\operatorname{Ric}(\Omega_t^M)-\Omega_t^M
 &=
 (2-\mu)\omega_{\rm tr}
 -
 A(t)\omega_{\rm tr}
 -
 e^{-t}s^\mu\sqrt{-1}\,\Theta\wedge\overline\Theta \\
 &=
 \bigl((2-\mu)-A(t)\bigr)\omega_{\rm tr}
 -
 e^{-t}s^\mu\sqrt{-1}\,\Theta\wedge\overline\Theta \\
 &= A'(t)\omega_{\rm tr} - e^{-t} s^\mu \sqrt{-1}\,\Theta\wedge\overline\Theta \\
 &= \partial_t\Omega_t^M.
\end{align*}
The initial value follows from \eqref{eq:Omega-flow-def}, because $A(0)=M$.

Finally, let $K\Subset S\setminus D$ be compact. On finitely many lifted
coordinate charts covering $K$, the functions $s$, $B$, and all their
derivatives are bounded. Hence
\begin{equation*}
 e^{-t}s^\mu\sqrt{-1}\,\Theta\wedge\overline\Theta
 \longrightarrow 0
 \qquad
 \text{in }\mathcal C^\infty(K).
\end{equation*}
Since $A(t)\to2-\mu$,
\begin{equation*}
 \Omega_t^M
 \longrightarrow
 (2-\mu)\omega_{\rm tr}
 \qquad
 \text{in }\mathcal C^\infty_{\rm loc}(S\setminus D).
\end{equation*}
The limiting form is positive in the $w$-direction and vanishes on
$\partial_z$. Hence it has rank one and kernel equal to the vertical
foliation.
\end{proof}

\begin{thebibliography}{99}
\bibitem[DOT03]{dloussky-oeljeklaus-toma} Georges Dloussky, Karl Oeljeklaus, and Matei Toma. \newblock Class ${\rm VII}_0$ surfaces with $b_2$ curves. \newblock \emph{Tohoku Math. J. (2)}, 55(2):283--309, 2003.
\bibitem[OT08]{oeljeklaus-toma-moduli} Karl Oeljeklaus and Matei Toma. \newblock Logarithmic moduli spaces for surfaces of class ${\rm VII}$. \newblock \emph{Math. Ann.}, 341(2):323--345, 2008.
\end{thebibliography}
\end{document}
